% Options for packages loaded elsewhere
\PassOptionsToPackage{unicode}{hyperref}
\PassOptionsToPackage{hyphens}{url}
\PassOptionsToPackage{dvipsnames,svgnames,x11names}{xcolor}
%
\documentclass[
  authoryear,
  preprint,
  onecolumn]{elsarticle}

\usepackage{amsmath,amssymb}
\usepackage{lmodern}
\usepackage{iftex}
\ifPDFTeX
  \usepackage[T1]{fontenc}
  \usepackage[utf8]{inputenc}
  \usepackage{textcomp} % provide euro and other symbols
\else % if luatex or xetex
  \usepackage{unicode-math}
  \defaultfontfeatures{Scale=MatchLowercase}
  \defaultfontfeatures[\rmfamily]{Ligatures=TeX,Scale=1}
\fi
% Use upquote if available, for straight quotes in verbatim environments
\IfFileExists{upquote.sty}{\usepackage{upquote}}{}
\IfFileExists{microtype.sty}{% use microtype if available
  \usepackage[]{microtype}
  \UseMicrotypeSet[protrusion]{basicmath} % disable protrusion for tt fonts
}{}
\makeatletter
\@ifundefined{KOMAClassName}{% if non-KOMA class
  \IfFileExists{parskip.sty}{%
    \usepackage{parskip}
  }{% else
    \setlength{\parindent}{0pt}
    \setlength{\parskip}{6pt plus 2pt minus 1pt}}
}{% if KOMA class
  \KOMAoptions{parskip=half}}
\makeatother
\usepackage{xcolor}
\setlength{\emergencystretch}{3em} % prevent overfull lines
\setcounter{secnumdepth}{5}
% Make \paragraph and \subparagraph free-standing
\ifx\paragraph\undefined\else
  \let\oldparagraph\paragraph
  \renewcommand{\paragraph}[1]{\oldparagraph{#1}\mbox{}}
\fi
\ifx\subparagraph\undefined\else
  \let\oldsubparagraph\subparagraph
  \renewcommand{\subparagraph}[1]{\oldsubparagraph{#1}\mbox{}}
\fi


\providecommand{\tightlist}{%
  \setlength{\itemsep}{0pt}\setlength{\parskip}{0pt}}\usepackage{longtable,booktabs,array}
\usepackage{calc} % for calculating minipage widths
% Correct order of tables after \paragraph or \subparagraph
\usepackage{etoolbox}
\makeatletter
\patchcmd\longtable{\par}{\if@noskipsec\mbox{}\fi\par}{}{}
\makeatother
% Allow footnotes in longtable head/foot
\IfFileExists{footnotehyper.sty}{\usepackage{footnotehyper}}{\usepackage{footnote}}
\makesavenoteenv{longtable}
\usepackage{graphicx}
\makeatletter
\def\maxwidth{\ifdim\Gin@nat@width>\linewidth\linewidth\else\Gin@nat@width\fi}
\def\maxheight{\ifdim\Gin@nat@height>\textheight\textheight\else\Gin@nat@height\fi}
\makeatother
% Scale images if necessary, so that they will not overflow the page
% margins by default, and it is still possible to overwrite the defaults
% using explicit options in \includegraphics[width, height, ...]{}
\setkeys{Gin}{width=\maxwidth,height=\maxheight,keepaspectratio}
% Set default figure placement to htbp
\makeatletter
\def\fps@figure{htbp}
\makeatother

\makeatletter
\makeatother
\makeatletter
\makeatother
\makeatletter
\@ifpackageloaded{caption}{}{\usepackage{caption}}
\AtBeginDocument{%
\ifdefined\contentsname
  \renewcommand*\contentsname{Table of contents}
\else
  \newcommand\contentsname{Table of contents}
\fi
\ifdefined\listfigurename
  \renewcommand*\listfigurename{List of Figures}
\else
  \newcommand\listfigurename{List of Figures}
\fi
\ifdefined\listtablename
  \renewcommand*\listtablename{List of Tables}
\else
  \newcommand\listtablename{List of Tables}
\fi
\ifdefined\figurename
  \renewcommand*\figurename{Figure}
\else
  \newcommand\figurename{Figure}
\fi
\ifdefined\tablename
  \renewcommand*\tablename{Table}
\else
  \newcommand\tablename{Table}
\fi
}
\@ifpackageloaded{float}{}{\usepackage{float}}
\floatstyle{ruled}
\@ifundefined{c@chapter}{\newfloat{codelisting}{h}{lop}}{\newfloat{codelisting}{h}{lop}[chapter]}
\floatname{codelisting}{Listing}
\newcommand*\listoflistings{\listof{codelisting}{List of Listings}}
\makeatother
\makeatletter
\@ifpackageloaded{caption}{}{\usepackage{caption}}
\@ifpackageloaded{subcaption}{}{\usepackage{subcaption}}
\makeatother
\makeatletter
\@ifpackageloaded{tcolorbox}{}{\usepackage[many]{tcolorbox}}
\makeatother
\makeatletter
\@ifundefined{shadecolor}{\definecolor{shadecolor}{rgb}{.97, .97, .97}}
\makeatother
\makeatletter
\makeatother
\ifLuaTeX
  \usepackage{selnolig}  % disable illegal ligatures
\fi
\usepackage[]{natbib}
\bibliographystyle{elsarticle-harv}
\IfFileExists{bookmark.sty}{\usepackage{bookmark}}{\usepackage{hyperref}}
\IfFileExists{xurl.sty}{\usepackage{xurl}}{} % add URL line breaks if available
\urlstyle{same} % disable monospaced font for URLs
\hypersetup{
  pdftitle={Simulated Land use and land cover change under scenarios of Ecological Infrastructure development in Switzerland between 2020-2060},
  pdfauthor={Benjamin Black; Adrienne Grêt-Regamey; Sven-Erik Rabe; Paula Mayer; Jan Streit},
  pdfkeywords={keyword1, keyword2},
  colorlinks=true,
  linkcolor={blue},
  filecolor={Maroon},
  citecolor={Blue},
  urlcolor={Blue},
  pdfcreator={LaTeX via pandoc}}

\setlength{\parindent}{6pt}
\begin{document}

\begin{frontmatter}
\title{Simulated Land use and land cover change under scenarios of
Ecological Infrastructure development in Switzerland between 2020-2060}
\author[1]{Benjamin Black%
\corref{cor1}%
}
 \ead{bblack@ethz.ch}
\author[1]{Adrienne Grêt-Regamey%
%
}
 \ead{gret@ethz.ch}
\author[1]{Sven-Erik Rabe%
%
}
 \ead{rabes@ethz.ch}
\author[1]{Paula Mayer%
%
}
 \ead{Mayerpa@student.ethz.ch}
\author[1]{Jan Streit%
%
}
 \ead{streitj@ethz.ch}

\affiliation[1]{organization={ETH Zürich, Planning of Landscape and
Urban Systems (PLUS) Institut für Raum- und Landschaftsentwicklung
(IRL)},addressline={HIL H 52.1, Stefano-Franscini-Platz
5},city={Zürich},postcode={8093},postcodesep={}}

\cortext[cor1]{Corresponding author}






\begin{abstract}
This is the abstract.
\end{abstract}





\begin{keyword}
    keyword1 \sep
    keyword2
\end{keyword}
\end{frontmatter}
    \ifdefined\Shaded\renewenvironment{Shaded}{\begin{tcolorbox}[breakable, interior hidden, sharp corners, frame hidden, boxrule=0pt, borderline west={3pt}{0pt}{shadecolor}, enhanced]}{\end{tcolorbox}}\fi

\defcitealias{sfso2021}{SFSO, 2021}

\hypertarget{introduction}{%
\section{Introduction}\label{introduction}}

The mountainous central European country of Switzerland, like the
majority of countries globally, exhibits declining biodiversity due to
multiple drivers such as urbanization, climatic change and
environmentally deleterious agricultural practices. To address
biodiversity loss, in 2012, the federal government of Switzerland
ratified the Swiss Biodiversity Strategy (SBS) \citet{swissbi2012} . One
of the goals of the SBS is the development of Switzerland's Ecological
Infrastructure (EI) as a ``network of protected and connected areas that
contribute to safeguarding the essential services of ecosystems for
society and economy \ldots{} and form the basis for a rich and resilient
biodiversity'' \citet{federalofficefortheenvironment}. Achieving this
goal will require consideration of not just how Switzerland's EI should
be planned under existing conditions but how it's development in space
and time could proceed in concurrence with changes in the boundary
conditions of the socio-ecological system within which it is situated.

In this regard, \citet{mayer2023} defined 5 narrative scenarios to frame
the development of EI in Switzerland between 2020 and 2060 alongside
expected demographic, economic and climatic changes in Switzerland. The
scenarios were produced through a participatory process with the aim of
envisioning nature-positive outcomes using the IPBES Nature futures
framework. As such the scenarios can characterized as both exploratory
and normative, encapsulating the development of a ``functioning'' EI
from different perspectives of societal-nature relationships (ibid).

However, whilst descriptive scenarios are a useful starting point for
planners and decision-makers, their legitimacy can be increased by
operationalising the narratives through computer-based simulation
modelling so that their spatio-temporal impacts on critical aspects of
the study system can be visualized and explored in detail
\citep{verburg2004}. Such operationalization of scenarios is especially
desirable when their realization would likely result in trade-offs that
must be negotiated by society. This is pertinent in the case of EI,
which, under the definition employed by the SBS, aims to balance both
ecosystem service ES) provision and biodiversity protection, which have
often been acknowledged as discordant objectives.

In terms of operationalizing the specific scenarios of EI development
from \citet{mayer2023}, the proximal linkages between EI as a `network
of areas' at the landscape level in Switzerland and it's intended
outcomes for biodiversity and ES are difficult to model directly.
Instead, it is logical to focus on a higher level of abstraction and
model an intermediary, or distal, interaction between the physical
expression of EI and these outcomes. A common component in the modelling
of both measures of biodiversity (e.g species distribution) and ES
provision \citep[\citet{gomes2021}]{newbold2015} is land use and land
cover (LULC). Furthermore LULC change (LULCC) trends are also a central
theme in the SBS concept of EI as well as the scenario narrative's of
\citet{mayer2023}. In particular, the role of protected areas and Other
Effective Conservation Measures (OECMs) as components of the EI as well
as other EI-relevant LULCC trends such as urban sprawl and the spatial
distribution of agricultural production. On this basis, the goal of this
research is to operationalize the scenario's of EI development produced
for Switzerland by \citet{mayer2023} by simulating their implications
for trends in LULCC in Switzerland between 2020 and 2060 using a
temporally dynamic Constrained Cellular Automata (CCA) model.

CCA models have enjoyed enduring popularity for the modelling of
scenarios of LULCC for the last X decades. LULCC-CCAs treat landscapes
as a pixelated grid of cells of individual LULC states with cells
changing state over time typically based on two processes. Firstly, the
spatial distribution of potential LULCC is determined through
calculation of cellular probabilities of change, often referred to as
transition potential (TP), to alternate LULC states. The means by which
this is achieved, differs between popular LULCC-CAs although it is
commonly based on the quantification of correlative relationships with
explanatory variables such as biophysical and socio-economic and
neighbourhood (focal LULC) predictors . Secondly, the amount of cellular
LULC changes that occur in a given time step is prescribed by values of
transition rates extrapolated from historic LULC data, with an
allocation algorithm selecting the required number of cells to change
based upon their transition probabilities .

The internal model process of LULCC-CCAs offer several opportunities to
operationalize aspects of scenarios within simulations that may be used
in isolation or in combination with each other. Firstly, the data
supplied in the form of dynamic variables used to predict cellular TP
values can be scenario specific. For example, in scenarios contrasting
differential levels of climatic change, the data provided may come from
simulations of different Relative Concentration Pathways . A second
approach is to alter predicted cellular TP values according to
proscribed criteria, This can be done directly: such as spatially
designated land management policies, for example reducing the potential
for transition to urban LULC states inside environmental protection
areas Or indirectly by altering the influence of variables within the
model used to calculate TPs .Finally, the influence of different
scenarios can be incorporated through specification of the transitions
rates\ldots. The influence of these three mechanisms within the
LULCC-CCA model are highlighted visually in Figure 1 below. To perform
scenario-based simulation of future LULCC in Switzerland between 2020
and 2060 we utilized a constrained cellular automata (CCA) modelling
approach combining the R statistical computing environment with the
environment modelling software Dinamica EGO .

\hypertarget{methods}{%
\section{Methods}\label{methods}}

\hypertarget{model-overview}{%
\subsection{Model overview}\label{model-overview}}

Figure~\ref{fig-process} highlights the primary processes taking place
within the LULCC-CCA model, the details of each these processes are
summarized in the subsequent sections although for brevity some material
has been presented as appendices.

\begin{figure}

{\centering \includegraphics{figures_tables/Fig1.jpg}

}

\caption{\label{fig-process}test}

\end{figure}

\hypertarget{land-use-land-cover-transition-identification}{%
\subsection{Land use land cover transition
identification}\label{land-use-land-cover-transition-identification}}

The process of identifying LULC transitions, i.e.~pixel changes from an
initial specific LULC class to a different final LULC class, was based
upon a preceding study by \citet{black_re-considering_2023}. We utilized
classified LULC datasets for four time periods from the Swiss area
statistics \citepalias{sfso2021}: 1981-1985; 1992-1997; 2004-2009;
2013-2018 (Note that date range is the extent of the data-collection
flying period, and in fact each product represents a single LULC map).
Given that transitions are identified between two time points these four
historic LULC maps meant that we had three LULC transition periods,
which will henceforth be referred to using the latter year of the flying
period as: 1985-1997; 1997-2009; and 2009-2018 respectively. Following
\citet{black_re-considering_2023} we aggregated the original 72 LULC
classes in the data to 10 classes: Alpine pastures; Closed forest;
Glacier; Grassland; Intensive agriculture; Open Forest; Shrubland;
Permanent crops; Settlement/urban/amenities; Static (see Appendix A for
a table of class aggregations)

\hypertarget{transition-potential-modelling}{%
\subsection{Transition potential
modelling}\label{transition-potential-modelling}}

\hypertarget{ca-parameter-calibration}{%
\subsection{CA parameter calibration}\label{ca-parameter-calibration}}

\hypertarget{scenario-development}{%
\subsection{Scenario development}\label{scenario-development}}

\hypertarget{scenario-quantification}{%
\subsection{Scenario quantification}\label{scenario-quantification}}

\hypertarget{transition-rates}{%
\subsubsection{Transition rates}\label{transition-rates}}

Price et al.~2017 use a similar process as we plan to calculate the
expected increases in the `most important' categories and then calculate
increates in the other categories from the remaining land.

Price et al.~2017 on calculating demand for urban area: ``To calculate a
total urban area demand per scenario, we used the modelled population
level for 2060 in combination with the minimum and maximum of current
per capita urban area demand data for the five cantons in our study
area.''

\hypertarget{validation}{%
\subsection{Validation}\label{validation}}

\hypertarget{simulation}{%
\subsection{Simulation}\label{simulation}}

\hypertarget{results}{%
\section{Results}\label{results}}

\hypertarget{discussion}{%
\section{Discussion}\label{discussion}}

\hypertarget{references}{%
\section{References}\label{references}}

\hypertarget{figures-and-tables}{%
\section{Figures and tables}\label{figures-and-tables}}

\hypertarget{references-1}{%
\section*{References}\label{references-1}}
\addcontentsline{toc}{section}{References}

ist


  \bibliography{bibliography.bib}


\end{document}
